\section{Previous Work}

\subsection{Alternative Methods to RAG}
A lot of recent work has been done on alternative methods to Retrieval-Augmented Generation (RAG) in an attempt to mitigate its weaknesses \citep{gao2023retrieval}. While RAG effectively grounds large language models in external knowledge, it fundamentally relies on vector similarity, which struggles with multi-hop reasoning, strict relational constraints, and explicit fact-checking \citep{edge2024identifying}. Consequently, approaches that integrate structured knowledge bases—particularly Knowledge Graphs (KGs)—with generative models have gained traction.

One prominent alternative is GraphRAG, which combines the unstructured text retrieval of standard RAG with the structural topological data of a knowledge graph \citep{edge2024identifying}. By indexing both raw text chunks and the entities/relations extracted from them, GraphRAG can answer complex, global questions over a dataset that standard vector retrieval fails to synthesize. However, GraphRAG still relies heavily on the LLM's latent semantic reasoning to bridge the gap between the retrieved graph structure and the final answer, meaning it does not inherently guarantee logical determinism.

Another approach is the translation of natural language to formal query languages (e.g., Text-to-SPARQL or Text-to-Cypher) \citep{kaufmann2014evaluating}. In these systems, the generative model does not answer the question directly; instead, it writes a database query, executes it against an ontology or graph database, and returns the deterministic result. While this guarantees factual accuracy under the Open World Assumption, it is extremely brittle. If the model hallucinates a relation name or misunderstands the ontology schema, the query fails entirely.

The GIG system builds upon these concepts but diverges by enforcing completely deterministic, schema-restricted data handoffs between agents rather than relying on query generation. By restricting the problem space to atomic questions (Section \ref{subsubsec: Atomic-question Assumption}) and using constrained generation to trace reasoning steps, GIG attempts to combine the strict determinism of a semantic web query with the fault tolerance of a multi-agent system.

\subsection{Small Language Models and Edge AI}
The push towards localized, privacy-preserving AI has driven significant research into Small Language Models (SLMs) and edge deployment. The Microsoft Phi series \citep{gunasekar2023textbooks}, which forms the backbone of the GIG system, demonstrated that high-quality, textbook-curated pretraining data allows sub-10-billion parameter models to rival much larger models in reasoning and instruction-following tasks.

When deploying these models on constrained hardware, research has increasingly focused on utilizing Neural Processing Units (NPUs) combined with aggressive quantization techniques. Gholami et al. \citep{gholami2021survey} outline how reducing model weights to INT4 precision drastically lowers memory bandwidth requirements, enabling real-time inference on edge devices without the power draw of dedicated GPUs. The GIG system explicitly tests these constraints by forcing an INT4-quantized SLM running locally on an NPU to perform complex, multi-stage ontological parsing—a task typically reserved for cloud-hosted LLMs.

\subsection{Neuro-Symbolic Task Decomposition}
The core agentic architecture of the GIG system is heavily influenced by research into neuro-symbolic multi-agent workflows. The Least-to-Most prompting framework \citep{zhou2022least} and the LAG (Logically Aligned Generation) framework demonstrate that language models are significantly more reliable when a complex problem is explicitly broken down into sequential, atomic sub-problems. 

Furthermore, constrained structured generation \citep{willard2023efficient} has emerged as an essential technique for stringing these sub-problems together. By forcing the neural model to output strict JSON schemas, systems can programmatically validate agent outputs before passing them to the next stage. The GIG system takes this to its logical conclusion: the neural agents are purely semantic processors (extracting entities, identifying relations), while the overall pipeline logic, knowledge storage, and final context mapping are handled entirely by deterministic, rule-based Python scripts. This strict separation of concerns is what enables the GIG system to maintain an auditable, hallucination-free trace from source text to final answer.
